% Part: second-order-logic
% Chapter: metatheory
% Section: undecidability-and-axiomatizability

\documentclass[../../../include/open-logic-section]{subfiles}

\begin{document}

\olfileid[pt]{sol}{met}{nax}

\olsection{A lógica de segunda ordem não é axiomatizável}


\begin{thm}
\ollabel{thm:sol-undecidable}
A lógica de segunda ordem é indecidível.
\end{thm}

\begin{proof}
Uma !!{sentence} de primeira ordem é válida na lógica de primeira ordem sse
é válida na lógica de segunda ordem, e a lógica de primeira ordem é
indecidível.
\end{proof}

\begin{thm}
\ollabel{cor:sol-not-axiomatizable}
Não existe um sistema de !!{derivation} correto e completo para a lógica de
segunda ordem.
\end{thm}

\begin{proof}
Seja $!A$ !!a{sentence} na linguagem da aritmética. $\Sat{N}{!A}$ sse
$\Th{PA^2} \Entails !A$. Seja $!P$ a conjunção dos nove axiomas de
~$\Th{PA^2}$. $\Th{PA^2} \Entails !A$ sse $\Entails !P \lif !A$, O.s.\@,
$\Sat{M}{!P \lif !A}$. Considere agora a !!{sentence}
$\lforall[z][\lforall[u][\lforall[u'][\lforall[u''][\lforall[L][(!P' \lif !A')]]]]]$
obtida substituindo $\Obj{0}$ por $z$, $\prime$ pela variável de função
unária $u$, $+$ e $\times$ pelas variáveis de função binária $u'$ e $u''$,
respectivamente, e $<$ pela variável de relação binária~$L$, e quantificando
universalmente. Ela é uma sentença válida da lógica pura de segunda ordem
sse a sentença original era válida, sse $\Th{PA^2} \Entails !A$, sse
$\Sat{N}{!A}$. Assim, se houvesse um sistema de prova correto e completo
para a lógica de segunda ordem, poderíamos usá-lo para definir uma
enumeração computável $f\colon \Nat \to \Sent[L_A]$ das !!{sentence}s
verdadeiras em~$\Struct{N}$. Essa função seria representável em~$\Th{Q}$
por alguma fórmula de primeira ordem~$!B_f(x, y)$. Então a
!!{formula}~$\lexists[x][!B_f(x, y)]$ definiria o conjunto das !!{sentence}s
de primeira ordem verdadeiras de~$\Struct{N}$, contradizendo o teorema de
Tarski.
\end{proof}



\end{document}

